\tikzset{
  lens/.style={draw, thick, ellipse, minimum width=0.2cm, minimum height=1.2cm},
  mirror/.style={draw, thick, rectangle, minimum width=0.1cm, minimum height=1cm, fill=gray!40},
  laser/.style={draw, thick, rectangle, minimum width=1cm, minimum height=0.4cm, fill=red!20},
}

%\usepackage{tikzoptics}

%\ProvidesPackage{tikzoptics}[2025/07/13 TikZ Optics Elements]
%\RequirePackage{tikz}
\usetikzlibrary{arrows.meta, positioning, calc, decorations.pathmorphing}
\usetikzlibrary{shapes.geometric}

% Couleurs par défaut
\definecolor{laserred}{RGB}{220,20,60}
\definecolor{glassblue}{RGB}{100,100,255}
\definecolor{mirrorgray}{gray}{0.6}

\usepgflibrary {shadings}
% Source laser (boîte)
\tikzset{
  laser source/.style={
    draw, thick, rectangle,
    minimum width=1.2cm, minimum height=0.7cm,
    shading=axis,
    left color=black, right color=black, middle color=colorTwo,
    %top color=black, middle color=white, bottom color=black,
    %outer color=black,inner color=white,
    shading angle=0
  }
}

% Style pour faisceau laser
\tikzset{
  laser beam/.style={
    thick,
    %draw=laserred,
    draw=colorSix, 
    %decoration={snake, amplitude=0.5mm, segment length=2mm},
    %decorate
  }
}
\tikzset{
  sond beam/.style={
    dashed,
    %draw=laserred,
    draw=colorSix, 
    %decoration={snake, amplitude=0.5mm, segment length=2mm},
    %decorate
  }
}


% Fibre optique (cercle)
\tikzset{
  fiber/.style={
    draw = none, thick, circle, minimum size=2.4cm,
    fill= none,
    color = colorSix, 
    path picture={
      \draw[->, thick]
        ([shift={(-0.2,0)}]path picture bounding box.center)
        arc[start angle=-90, end angle=300, radius=0.5cm];
      \draw[->, thick]
        ([shift={(0.2,0)}]path picture bounding box.center)
        arc[start angle=-90, end angle=300, radius=0.5cm];
    }
  }
}


% Lentille convergente (symbole)
\tikzset{
  lens/.style={
    draw = none , thick, fill= none ,
    minimum width=0.2cm, minimum height=1.4cm,
    shape=ellipse,
    path picture={
      \draw[ thick , shading=axis , color = colorFour , left color=colorFour, right color=colorFour, middle color=colorTwo, shading angle=0]
        ([shift={(-0.1,0)}]path picture bounding box.center)
        arc[start angle=180, end angle=170, radius=3.5cm] --++ (0.1, 0 ) 
        arc[start angle=10, end angle=-10, radius=3.5cm] --++ (-0.1 , 0 )
        arc[start angle=-170, end angle=-180, radius=3.5cm] 
       ;
       %\drawgrid{-3}{3}{-1}{3};
    }
  }
}

\usetikzlibrary{patterns.meta}
 %Miroir plan (rectangle vertical incliné)
%\tikzset{
%  mirror/.style={
%    draw = none , thick, fill= none ,
%    minimum width=0.4cm, minimum height=1.4cm,
%    shape=rectangle, rotate=0,
%    path picture={
%      \draw[thick , line width=0.2ex,
%      	color = colorFour , left color=colorFour, right color=white, shading angle=90
%      	]
%        ([shift={(-0.,-0.6)}]path picture bounding box.center)
%        rectangle (0.1,0.6);
%      \fill[
%  		pattern=north west lines,
%  		pattern color=colorFour,
%  		%pattern line width = 0.3pt,     % finesse des traits
%		%pattern scale = 0.5             % plus petit motif
%  		]
%        ([shift={(0.1,-0.6)}]path picture bounding box.center)
%        rectangle ([shift={(0.15,0.6)}]path picture bounding box.center);
%      %\drawgrid{-3}{3}{-1}{3};
%    } 
%  }
%}

%\tikzset{
%  mirror/.style={
%    draw=none, thick, fill=none,
%    minimum width=0.4cm, minimum height=1.4cm,
%    shape=rectangle, rotate=0,
%    path picture={
%      % Bande colorée
%      \draw[thick, line width=0.2ex,
%        color=colorFour, left color=colorFour, right color=white, shading angle=90]
%        ([shift={(-0.,-0.6)}]path picture bounding box.center)
%        rectangle (0.1,0.6);
%
%      % Zone hachurée avec un clip
%      \scope
%        \clip
%          ([shift={(0.1,-0.6)}]path picture bounding box.center)
%          rectangle
%          ([shift={(0.15,0.6)}]path picture bounding box.center);
%
%        % Hachures à 45°
%        \foreach \t in {-1.0,-0.8,...,1.0} {
%          \draw[color=colorFour, line width=0.2pt]
%            ({\t-1},-1) -- ({\t+1},1);
%        }
%      \endscope
%    }
%  }
%}

\tikzset{
  mirror/.style={
    draw=none, line width=1.0pt, fill=none,
    minimum width=0.4cm, minimum height=1.4cm,
    shape=rectangle, rotate=0,
    path picture={
      % Bande colorée
      \draw[thick, line width=0.2ex,
        color=colorFour, left color=colorFour, right color=colorTwo, shading angle=90]
        ([shift={(-0.0,-0.6)}]path picture bounding box.center)
        rectangle (0.1,0.6)
        ([shift={(0.12,-0.615)}]path picture bounding box.center) -- ([shift={(0.12,0.615)}]path picture bounding box.center)
      ;        
      % Hachures à 45° (pas de ligne vide ni de \par)
      \scope
        \clip
          ([shift={(0.1,-0.6)}]path picture bounding box.center)
          rectangle
          ([shift={(0.2,0.6)}]path picture bounding box.center);
        % Hachures à 45°
        \foreach \t in {-1.0,-0.9,...,1.0} {
          \draw[color=colorFour, line width=1.0pt]
            ({\t-1},-1) -- ({\t+1},1);
        }
      \endscope
    }
  }
}


\tikzset{
  sep/.style={
    draw=none, line width=1.0pt, fill=none,
    minimum width=0.4cm, minimum height=1.4cm,
    shape=rectangle, rotate=0,
    path picture={
      % Bande colorée
      \draw[thick, line width=0.2ex,
        rounded corners=0.2ex , color = colorFour , left color=colorFour, right color=colorFour , middle color=colorTwo, shading angle=90]
        ([shift={(-0.0,-0.6)}]path picture bounding box.center)
        rectangle (0.15,0.6)
        %([shift={(0.12,-0.615)}]path picture bounding box.center) -- ([shift={(0.12,0.615)}]path picture bounding box.center)
      ; 
   }
   } 
}



% Lame demi-onde (rectangle avec double flèche)
\tikzset{
  halfwave plate/.style={
    draw , thick, rectangle,
    minimum width=0.2cm, minimum height=1.2cm,
    fill = none,
    line width=0.2ex , rounded corners=0.2ex , color = colorFive , left color=colorFive, right color=colorFive , middle color=colorTwo , shading angle=90,
%    path picture={
%      \draw[ line width=0.2ex , rounded corners=0.5ex , color = colorFive , left color=colorFive, right color=colorFive , middle color=white , shading angle=90]
%        ([xshift=0.1cm, yshift=0.1cm]path picture bounding box.south west)
%          rectangle
%        ([xshift=-0.1cm, yshift=-0.1cm]path picture bounding box.north east);
%    }
  }
}



% TA
\tikzset{
  TA/.style={
    draw = none , thick, rectangle,
    minimum width=1.2cm, minimum height=0.6cm,
    fill = none,
    line width=0.2ex , rounded corners=0.0ex , color = colorOne , left color=colorOne, right color=colorOne , middle color=colorTwo , shading angle=0,
    rotate = 0, 
    path picture={
%       \draw[ line width=0.2ex , rounded corners=0.0ex , color = colorFive , left color=colorFive, right color=colorFive , middle color=colorSix , shading angle=90]
%        ([xshift=0.0cm, yshift=0.01cm]path picture bounding box.west)
%          -- ([xshift=0.0cm, yshift=0.01cm]path picture bounding box.center) --
%        ([xshift=-0.0cm, yshift=0.2cm]path picture bounding box.east) -- ([xshift=-0.0cm, yshift=-0.2cm]path picture bounding box.east) -- ([xshift=0.0cm, yshift=-0.01cm]path picture bounding box.center) --  ([xshift=0.0cm, yshift=-0.01cm]path picture bounding box.west)
%       ;
       \draw[ line width=0.2ex , rounded corners=0.0ex , color = colorFive , left color=colorFive, right color=colorFive , middle color=colorSix , shading angle=90]
        ([xshift=0.0cm, yshift=0.2cm]path picture bounding box.west)
          -- ([xshift=0.0cm, yshift=0.01cm]path picture bounding box.center) --
        ([xshift=-0.0cm, yshift=0.01cm]path picture bounding box.east) -- ([xshift=-0.0cm, yshift=-0.01cm]path picture bounding box.east) -- ([xshift=0.0cm, yshift=-0.01cm]path picture bounding box.center) --  ([xshift=0.0cm, yshift=-0.2cm]path picture bounding box.west)
       ;
    }
  }
}

%      \draw[ line width=0.2ex , rounded corners=0.0ex , color = colorOne , left color=colorOne, right color=colorOne , middle color=colorTwo , shading angle=90]
%        ([xshift=0.1cm, yshift=0.1cm]path picture bounding box.south west)
%          rectangle
%        ([xshift=-0.1cm, yshift=-0.1cm]path picture bounding box.north east)
%       ;

% Lentille cylindrique (ellipse aplatie)
\tikzset{
  cylindrical lens/.style={
    draw = none , thick, fill= none ,
    minimum width=0.2cm, minimum height=1.4cm,
    shape=ellipse,
    path picture={
      \draw[ thick , line width=0.2ex , shading=axis , color = colorFour , left color=colorFour, right color=colorFour, middle color=white, shading angle=0]
        ([shift={(-0.1,0)}]path picture bounding box.center)
        arc[start angle=180, end angle=170, radius=3.5cm] --++ (0.1, 0 ) 
        arc[start angle=10, end angle=-10, radius=3.5cm] --++ (-0.1 , 0 )
        arc[start angle=-170, end angle=-180, radius=3.5cm] 
       ;
       %\drawgrid{-3}{3}{-1}{3};
    }
  }
}

% pignol
\tikzset{
  pignol/.style={
    draw = none , thick, fill= none ,
    minimum width=0.2cm, minimum height=1.4cm,
    shape=ellipse,
    path picture={
      \draw[ thick , line width=0.2ex , color = colorOne]
        ([shift={(0,0)}]path picture bounding box.north) -- ([shift={(0,0.1cm)}]path picture bounding box.center) 
       ;
      \draw[ thick , line width=0.2ex ,  , color = colorOne]
        ([shift={(0,0)}]path picture bounding box.south) -- ([shift={(0,-0.1cm)}]path picture bounding box.center) 
       ;
       %\drawgrid{-3}{3}{-1}{3};
    }
  }
}

% Isolateur optique (boîte avec flèche croisée)
\tikzset{
  optical isolator/.style={
    draw, thick, rectangle, minimum width=1cm, minimum height=0.7cm,
    fill=none ,
    line width=0.2ex , rounded corners=0.0ex , color = colorOne , left color=colorOne, right color=colorOne , middle color=colorTwo , shading angle=90,
    path picture={
      \draw[<-, thick] 
        ([shift={(-0.3,0.2)}]path picture bounding box.center)
        -- ([shift={(0.3,0.2)}]path picture bounding box.center);
        \draw[<-, thick] 
        ([shift={(-0.3,-0.2)}]path picture bounding box.center)
        -- ([shift={(0.3,-0.2)}]path picture bounding box.center);
%      \draw[<-, thick, red] 
%        ([shift={(0.3,-0.3)}]path picture bounding box.center)
%        -- ([shift={(-0.3,-0.3)}]path picture bounding box.center);
    }
  }
}

% Modulateur AOM (boîte avec flèche diagonale)
\tikzset{
  AOM/.style={
    draw, thick, minimum width=0.4cm, minimum height=1cm,
    shape=rectangle, fill=none,
    line width=0.2ex , rounded corners=0.0ex , color = colorOne , left color=colorOne, right color=colorOne , middle color=colorTwo , shading angle=90,
    path picture={
      \fill[pattern=horizontal lines,
  		pattern color=colorOne] 
        ([shift={(0.1,0.1)}]path picture bounding box.south west)
        rectangle ([shift={(-0.1,-0.1)}]path picture bounding box.north east);
    }
  }
}


% Polariseur (rectangle fin + symbole)
\tikzset{
  polarizer/.style={
    draw, thick, minimum width=0.3cm, minimum height=1.4cm,
    fill=gray!30,
    path picture={
      \draw[path picture bounding box]
        ([xshift=-2pt]path picture bounding box.south)
        -- ([xshift=2pt]path picture bounding box.north);
      
    }
  }
}

% Collimateur
\tikzset{
  collimator/.style={
    draw, thick, trapezium, trapezium angle=60,
    minimum height=1cm, fill=blue!10
  }
}


% Miroir dichroïque
\tikzset{
  dichroic/.style={
    draw, thick, fill=yellow!30, rectangle,
    minimum width=0.2cm, minimum height=1.4cm, rotate=45
  }
}

% Détecteur photodiode
\tikzset{
  photodiode/.style={
    draw, thick, regular polygon, regular polygon sides=3, fill=orange!30,
    minimum size=1.2cm, rotate=270
  }
}


% Pinhole (trou)
\tikzset{
  pinhole/.style={
    draw, thick, circle, minimum size=0.3cm, fill=black
  }
}


%DMD
\tikzset{
  DMD/.style={
    draw=none, line width=1.0pt, fill=none,
    minimum width=0.5cm, minimum height=1.5cm,
    shape=rectangle, rotate=0,
    path picture={
      % Bande colorée
      \draw[thick, line width=0.2ex,
        color=colorPrune, left color=colorPrune, right color=white, shading angle=90]
        ([shift={(-0.,-0.7)}]path picture bounding box.center)
        rectangle (0.1,0.7);
      % Hachures à 45° (pas de ligne vide ni de \par)
      \scope
        %\clip
%        \draw[thick]
%          ([shift={(0.1,-0.6)}]path picture bounding box.center)
%          rectangle
%          ([shift={(0.15,0.6)}]path picture bounding box.center);
        % Hachures à 45°
        \def\listetrais{-1.0,-0.9,...,-0.2, 0.1 ,0.2,...,1.0 }
        \def\traity{0.05}
        \def\alpha{12}
        \foreach \t in \listetrais {
		\draw[shift={(1,\t)}, rotate=12, color=colorPrune, line width=0.5pt] 
          ({0.18 - \traity*sin(\alpha)},{\t- \traity*cos(\alpha)}) -- ({0.18 + \traity*sin(\alpha)},{\t+ \traity*cos(\alpha)});
		}
		\def\listetrais{-0.2,-0.1,...,0}
		\foreach \t in \listetrais {
		\draw[shift={(1,\t)}, rotate=12, color=colorPrune, line width=0.5pt] 
          ({0.18 - \traity*sin(-\alpha)},{\t- \traity*cos(-\alpha)}) -- ({0.18 + \traity*sin(-\alpha)},{\t+ \traity*cos(-\alpha)});
		}
      \endscope
    }
  }
}

%lens deme
\tikzset{
  lensd/.style={
    draw = none , thick, fill= none ,
    minimum width=0.2cm, minimum height=1.4cm,
    shape=ellipse,
    path picture={
      \draw[ thick , line width=0.2ex , shading=axis , color = colorFour , left color=colorFour, right color=colorFour, middle color=white, shading angle=0]
        ([shift={(0.0,0.6)}]path picture bounding box.center)
         --++ (0.05, 0 ) 
        arc[start angle=10, end angle=-10, radius=3.5cm] --++ (-0.05 , 0 ) -- cycle
       ;
       %\drawgrid{-3}{3}{-1}{3};
    }
  }
}


%lens deme
\tikzset{
  lensbess/.style={
    draw = none , thick, fill= none ,
    minimum width=1.2cm, minimum height=2.4cm,
    shape=ellipse,
    path picture={
      	\node[lens] (L) {}; 
      	\node[lensd ,  right=-2mm of L]  {};
      	\node[lensd , rotate = 180,  right=-1.5mm of L]  {};
       %\drawgrid{-3}{3}{-1}{3};
    }
  }
}

%lens deme
\tikzset{
  atome/.style={
    draw = red , thick, fill= none ,
    %minimum width=0.1cm, minimum height=0.1cm,
    inner sep=0pt,    % pas de padding interne
    outer sep=0pt,    % pas de marge externe
    %minimum width=0pt, 
    %minimum height=0pt,
    shape=ellipse,
    path picture={
      	\draw[ thick , line width=0.2ex , shading=axis , color = colorFour , left color=colorFour, right color=colorFour, middle color=colorFour, shading angle=0]
      	([shift={(0.0,0)}]path picture bounding box.center) circle (0.3);
       %\drawgrid{-3}{3}{-1}{3};
    }
  }
}


\tikzset{
  CCD/.style={
    draw, thick, rectangle,
    minimum width=1.2cm, minimum height=0.7cm,
    shading=axis,
    left color=colorOne, right color=colorOne, middle color=colorTwo,
    %top color=black, middle color=white, bottom color=black,
    %outer color=black,inner color=white,
    shading angle=0
  }
}

